\SourceLeaf{246}%
also im \textsc{Lobatschewsky}schen Werke nicht gefunden, aber die Entwickelung
ist auf anderm Wege gemacht, als ich selbst eingeschlagen habe, und zwar von
\textsc{Lobatschewsky} auf eine meisterhafte Art in ächt geometrischem Geiste.
Ich glaube Sie auf das Buch aufmerksam machen zu müssen, welches Ihnen gewiss
ganz exquisiten Genuss gewähren wird. \mbox{.\ .\ .\ .\ .\ .}

\par\bigskip\Needspace{18\baselineskip}

\begin{center}
\textsc{Gauss} an W. \textsc{Struve}.\quad Göttingen, 11. December 1846.
\end{center}

\mbox{.\ .\ .\ .\ .\ .\ .} Gleichermassen bin ich für die übrigen Zusendungen zu dem
verbindlichsten Danke verpflichtet; für die russischen Sachen von
\textsc{Lobatschewsky} wahrscheinlich zunächst Ihrem Herrn Sohne, gegen den
ich vor einigen Jahren bei seinem Hiersein meinen Wunsch ausgesprochen hatte;
ich lasse mich seinem freundlichen Andenken angelegentlich empfehlen. Mit
meiner russischen Sprachkenntniss werde ich wohl etwas zurückgekommen sein,
da ich seit länger als einem Jahre nicht dazu habe kommen können, auch nur
einen russischen Buchstaben anzusehen, ich hoffe jedoch in der ersten freien
Zeit das Versäumte schnell nachzuholen, und werde dann der Lecture jener
interessanten Aufsätze meine besondere Aufmerksamkeit widmen. Die kleine
deutsche Schrift von \textsc{Lobatschewsky} besass ich schon vorher
selbst. \mbox{.\ .\ .\ .\ .\ .\ .}


